\hypertarget{into-the-machine}{%
\chapter{1. Into the machine}\label{into-the-machine}}

The design volume set out what a kennel is: a reference monitor at the
user level, a mediator every access passes through that the workload
cannot bypass or tamper with, built from two limbs. One limb is
construction, where a resource is simply absent from the kennel's world.
The other is interposition, where a resource that cannot be removed is
reachable only through a transaction the monitor authorises. None of
that named a kernel.

A modern Linux kernel offers a set of facilities that, taken one at a
time, each do a small part of this: hide a class of resource, mediate
access to another, hold a privilege for a moment. No single one of them
is a sandbox. The confinement is the weave: the facilities combined,
resource class by resource class, so that each class is either
constructed away or placed behind a mediated transaction. What follows
names the threads and shows how they cross; the later chapters pick up
each in turn.

\hypertarget{what-counts-as-a-control}{%
\section{1.1 What counts as a control}\label{what-counts-as-a-control}}

Not every restriction the kernel can express is a control in the sense
the reference monitor needs. The test is narrow and the design holds to
it: a control counts only if it is enforced by the kernel and keyed to
process identity, such that the confined workload cannot forge or bypass
it from inside. A rule the workload sidesteps by changing its source
address, rewriting its environment, or choosing a different path to open
is a speed bump, not a control, and the design does not count it as one.

That test is the tamperproof property of the reference monitor
(\texttt{reference-monitor}) restated in kernel terms, and it does the
sorting. It is why the facilities below are the ones that carry the
limbs and other plausible-looking mechanisms are not: a check the
workload runs against itself, a convention it is asked to honour, a
filter it can route around all fail the test by construction. It is also
why some facilities appear only as fallbacks. Where the first-choice
mechanism is unavailable on an older kernel, a coarser one may carry the
class at reduced fidelity, and where even that is missing the policy
does not run at all rather than pretend (\texttt{refuse-to-start}).

\hypertarget{the-construction-facilities}{%
\section{1.2 The construction
facilities}\label{the-construction-facilities}}

The construction limb is built from namespaces, and its character is
absence: what was never mapped or never bound into the kennel is not
there to be reached, mediated or not (\texttt{construction-by-absence}).

The user namespace is the foundation the rest stands on. A kennel runs
in its own, created as the operator so that the namespace is
operator-owned, with precise identity maps rather than a subuid range:
the host root line and the operator's own line, plus one per granted
group. That gives the kennel a real uid 0 and \texttt{CAP\_SYS\_ADMIN}
inside the namespace and nowhere else, which is what permits the
\texttt{mount} and \texttt{pivot\_root} that build the view. The mapping
of host uid 0 is the one privileged act in the whole construction, and
it is the reason a kennel is built by a privileged step rather than by
the workload itself (T5.4). The mount namespace carries the constructed
view: the kennel sees the filesystem its policy binds in and a root
pivoted away from the host's, so a path that was not granted is not
merely denied but absent. The PID namespace makes the kennel's first
process PID 1 with a fresh \texttt{/proc}, so the host's processes are
not visible to enumerate, and the IPC and network namespaces remove the
System V and network surfaces the same way. Each namespace takes a class
of resource out of the kennel's world rather than guarding it, which is
the cheaper and stronger move wherever it can be made.

\hypertarget{the-interposition-facilities}{%
\section{1.3 The interposition
facilities}\label{the-interposition-facilities}}

Some resources cannot be removed, because the kennel exists to use them:
it must execute some binaries, reach some paths, open some network
destinations. For these the design does not remove the class but places
it behind a mediated transaction (\texttt{interpose-as-transaction}).
Four kernel facilities carry that limb. Three are the standard sandbox
vocabulary and are described here; the fourth, the one that does not
look like it belongs, gets its own section after.

Landlock is the principal one. It applies a kernel-enforced,
unforgeable, identity-keyed access policy to the filesystem (read,
write, and with \texttt{FS\_EXECUTE} a proper execute semantics), to
network ports, and on ABI 6 to the scoping of abstract AF\_UNIX sockets
and signals. It is the mechanism that lets a kennel hold a path it may
read but not write, or execute one binary and not another, with the
decision made in the kernel and not in code the workload can reach.
Seccomp carries the syscall surface that Landlock does not express: the
filter that refuses \texttt{ptrace}, the \texttt{TIOCSTI}
terminal-injection guard, and the \texttt{connect()} restriction that
stands in for abstract-socket scoping on kernels below ABI 6. The cgroup
BPF programs carry network egress: the \texttt{inet4\_connect} and
\texttt{inet6\_connect} hooks enforce the address allowlist in the
kernel, on the connecting process, with no path around them, and this is
the egress floor present in every network mode but the fully isolated
one. Those three are the interposition vocabulary a sandbox reader
already knows. The fourth is not.

\hypertarget{binder}{%
\section{1.4 Binder}\label{binder}}

The fourth interposition facility will stop a reader who knows the
others cold, because it comes from Android: binder, exposed on Linux as
binderfs. Its presence is not exotic taste, and the reason is worth
setting out, because binder carries the whole local-services and
inter-kennel story in the chapters ahead.

The problem it solves is one a socket cannot. Several of the things a
kennel must talk to, the session bus or a sibling service kennel, speak
a protocol that grants too much to forward raw while the application
expects a conforming endpoint at a familiar path. Forwarding the raw
socket makes the grant auditable once, at connect time, and then
surrenders every call made across it. What the design needs is a
chokepoint at the call rather than the connection: one point where every
transaction a kennel makes with a service is mediated and audited, with
kenneld the decision point for each, and where a kennel cannot reach a
service it was not granted. The same primitive has to carry
communication between kennels, a client kennel calling a standing
service kennel, with neither holding ambient access to the other.

binderfs supplies all of that off the shelf. It is the kernel's
per-mount-namespace Binder filesystem, stable well under the 6.10 floor
the kennel already requires for Landlock. It gives synchronous
call-and-return transactions with kernel-managed identity and no
application framing; a service registry brokered by a well-known node
that kenneld occupies as context manager; references that are opaque
kernel objects a process either holds or does not, with no path to
enumerate and no abstract name to probe, so a reference passes the test
of §1.1 by construction; exact kernel-driven death notification; and
per-instance isolation, each binderfs mount a fully independent instance
the way devpts and tmpfs are, so two kennels' buses share no state. The
property that turns it from a curiosity into a fit is that binderfs is
\texttt{FS\_USERNS\_MOUNT}: it mounts unprivileged inside the kennel's
own user namespace, so the inter-namespace gateway is kernel-mediated
and unforgeable and yet cost no host capability to stand up.

What the design declines is the part of Binder that is hard. Android
passes node references between processes and maintains the capability
graph that results, which is most of what makes that implementation
subtle. The kennel uses none of it: references are issued by kenneld as
context manager and do not transfer between kennels, which keeps the
properties that matter, unforgeable references, synchronous calls, death
notification, per-instance isolation, and drops the bookkeeping that
does not earn its place.

That leaves the real question, which is why adopt an IPC system from
another world at all rather than write a small bus over socketpairs. The
answer is the supply discipline (\texttt{dont-roll-your-own}). A
homebrew bus would reimplement transaction routing, reference
management, death notification, and instance isolation that binderfs
already provides and that Android has run at a scale no in-tree bus will
see, and the kennel already speaks ioctl to the kernel for Landlock and
BPF, so binder's ioctl ABI is the same shape of code rather than a new
dependency. The unfamiliar primitive is the conservative choice, not the
adventurous one. How kenneld holds the context-manager node, how the
facades transact across the boundary, and the wire format itself come
later; here it is enough that the strangest thread in the weave is there
because it is the one that fits.

\hypertarget{privilege-and-identity}{%
\section{1.5 Privilege and identity}\label{privilege-and-identity}}

The third group of facilities is not a limb but the handling of
privilege itself, and the design's posture is to hold it briefly and in
one place. File capabilities let the one privileged binary carry exactly
the capabilities its construction step needs and no general root, so the
privilege that builds a kennel is a named, narrow set rather than
ambient authority (\texttt{rule-of-1}). \texttt{PR\_SET\_NO\_NEW\_PRIVS}
is set unconditionally before the workload runs: it is the precondition
seccomp requires, and it is the kernel's bar on a process regaining
privilege through a later \texttt{exec}, so a workload cannot climb back
up through a setuid binary it finds. The uid and gid maps and the
\texttt{capset} that drops the bounding set are what take the
constructed kennel from its transient uid 0 down to the operator's
identity with nothing attached to it the policy did not grant. This is
\texttt{split-the-uid} in mechanism: the kennel keeps an identity and is
stripped of the authority that identity carried on a traditional system.
How these are sequenced, which binary holds which capability, and for
how long, is the subject of the next chapter.

\hypertarget{the-weave}{%
\section{1.6 The weave}\label{the-weave}}

No single facility is sufficient, and the design says so plainly: the
mechanisms are used in combination, resource class by resource class.
The map is the architecture. Filesystem access is a mount namespace for
the view and Landlock for the access inside it. Execution is Landlock
\texttt{FS\_EXECUTE} and \texttt{PR\_SET\_NO\_NEW\_PRIVS} together. A
network address is the cgroup BPF connect hooks with no fallback,
because nothing coarser meets the test. Abstract AF\_UNIX is Landlock
scoping on ABI 6 and a seccomp \texttt{connect()} filter below it.
Process visibility is a mount namespace with \texttt{hidepid} and a PID
namespace for the stronger cut. Most classes are carried by a primary
mechanism with a fallback named beside it, and a few are carried by one
mechanism with no fallback at all, which means the kernel must provide
it or the kennel does not start.

The weave is the reason the confinement degrades gracefully in fidelity
and never silently in safety. Pull a thread on an older kernel and the
class it carried often retains a coarser fallback, at a fidelity the
policy can see and a residual the catalogue records. Pull a thread that
has no fallback and the policy refuses to apply rather than run a kennel
that looks confined and is not. Complete mediation is this map holding
for every class at once; the verifiable property is that where it cannot
hold, the system declines rather than degrades
(\texttt{refuse-to-start}).

\hypertarget{the-kernel-floor}{%
\section{1.7 The kernel floor}\label{the-kernel-floor}}

Because the weave depends on specific kernel facilities, a kennel
depends on a specific minimum kernel, and the design states the floor
rather than discovering it at runtime. Landlock is the load-bearing
thread, and its versions set the floor: the filesystem ACL arrives at
5.13, network ports at 6.7, \texttt{FS\_EXECUTE} at 6.10, and the ABI 6
scoping that natively isolates abstract AF\_UNIX and signals at 6.12.
The project floor is 6.10, ABI 5; on 6.12 and above the scoping
fallbacks fall away and the isolation is native. The cgroup v2
hierarchy, the cgroup BPF connect hooks, and the namespaces are all
long-standing and effectively universal on a modern system. At policy
load the kennel checks the features its policy requires against the
kernel it is on, and where a required feature is missing it refuses to
apply the policy and names what is absent rather than failing obscurely
or proceeding without the control.

One requirement is a deployment property rather than a kernel version.
The construction rests on an unprivileged user namespace, and some
hardened distributions restrict that by default. Where the restriction
is in force the namespace can be created but holds no capabilities
inside it, so the first map write fails and no kennel can be built. The
supported posture is an AppArmor profile that grants user-namespace
creation to the daemon binary alone, restoring the one capability
without relaxing the host-wide restriction. It is a one-time install
step, the complement to the file capabilities the privileged binary
carries, and where the restriction is absent no profile is needed.

\hypertarget{the-same-weave-elsewhere}{%
\section{1.8 The same weave elsewhere}\label{the-same-weave-elsewhere}}

Naming these mechanisms does not make the design Linux; it makes this
volume Linux. The reference monitor and its two limbs are
platform-neutral, and the porting assessments weave the same limbs from
another kernel's facilities, graded against the same test from §1.1. On
macOS the only viable enforcement layer is Seatbelt, undocumented and
version-volatile, and it grades inferior to Linux on recon-resistance
and on same-uid loopback isolation, with the shortfall recorded as a
residual rather than hidden. On FreeBSD, jails with VNET and Capsicum
grade equivalent or superior on most controls, with the private network
stack turning the egress path into a trusted plumbing problem the port
must solve deliberately. The threads differ; the weave is the same
shape, and a port that graded equal across the board would be hiding
exactly the residuals an honest assessment exists to surface.

With the toolkit and the weave in hand, the rest of this part picks up
the threads one at a time. The first is the set of processes that do the
weaving and the single one among them that holds privilege.
